\documentclass[12pt]{article}
\usepackage[utf8]{inputenc}
\usepackage[T1]{fontenc}
\usepackage{hyperref}
\title{Suy\'a Language Resources for K-12 Curriculum Development}
\author{Compiled by Claude Sonnet 5.5 (claude-sonnet-5-5) for Andr\'as Kornai}
\date{}
\begin{document}
\maketitle

\section{Introduction}
Suy\'a (Kĩs\^edj\^e; ISO 639-3 suy; Glottocode suya1243) is a Northern J\^e language of Mato Grosso, Brazil, spoken in the Xingu Indigenous Park region. About 350 speakers are reported. Searches found little teaching material. This report lists what was found and does not pad.

\section{Primary Educational Resources}
No primer, textbook or school curriculum was found in the search results. A 2022 article in \emph{Tra\c{c}os de Linguagem} (in Portuguese) examines language policies and the language in the school context of the Khikatxi village: \url{https://periodicos.unemat.br/index.php/tracos/article/download/11220/7980/41067}

\section{Linguistic and Reference Materials}
\begin{itemize}
\item Wikipedia entry on Suy\'a: \url{https://en.wikipedia.org/wiki/Suy%C3%A1_language}
\item Wikipedia entry on Kĩs\^edj\^e: \url{https://en.wikipedia.org/wiki/K%C4%A9s%C3%AAdj%C3%AA_language}
\item Signum (UEL) articles that discuss the grammatical description, citing Guedes (1993) and Santos (1997, doctoral dissertation, UFSC): \url{https://ojs.uel.br/revistas/uel/index.php/signum/article/view/47304}
\item Endangered Languages Project bibliography: \url{https://endangeredlanguages.com/lang/2725/bibliography?hl=de}
\end{itemize}
No dictionary was found.

\section{Religious and Community Resources}
None found in the search results.

\section{Digital and Audio Resources}
No audio, video or corpus resources were found.

\section{Curriculum Development Context}
Documentation exists mainly as academic grammar and phonology. Curriculum development should be done with the Kĩs\^edj\^e community and its school staff. Further searching under the name Kĩs\^edj\^e in Portuguese may turn up more.
\end{document}